\documentclass[12pt]{article}
\usepackage[utf8]{inputenc}
\usepackage[T1]{fontenc}
\usepackage[margin=1in]{geometry}
\usepackage{setspace}
\usepackage{xcolor}
\usepackage[colorlinks=true, urlcolor=blue!60!black, linkcolor=blue!60!black]{hyperref}
\setlength{\parindent}{0em}
\setlength{\parskip}{0.45em}
\renewcommand{\baselinestretch}{1.0}
\begin{document}
\begin{center}
{\Large \textit{Nixon v. United States}, 506 U.S. 224 (1993)\par}
\vspace{0.3em}
{\normalsize Case Brief --- One-Page Edition\par}
\end{center}

\textbf{Memory Jogger.} The Senate's ``sole Power to try all Impeachments'' makes challenges to its trial procedures nonjusticiable (\href{https://www.courtlistener.com/opinion/112803/nixon-v-united-states/}{CourtListener}).

\textbf{Facts.} Chief Judge Walter L. Nixon, Jr., of the Southern District of Mississippi was convicted of lying to a grand jury investigating whether he took a gratuity to ask a prosecutor to drop a case. He refused to resign, kept drawing his salary in prison, and was impeached. Under Senate Rule XI, a committee took testimony from ten witnesses and reported; the full Senate heard argument and convicted him by more than two-thirds.

\textbf{Procedural History.} Nixon sued to void his removal; the district court and D.C. Circuit held the claim a political question.

\textbf{Votes.} Affirmed 9--0. Rehnquist wrote for six Justices; Stevens concurred; White (with Blackmun) and Souter concurred in the judgment.

\textbf{Holding.} The Clause textually commits trial procedure to the Senate, and ``try'' lacks judicially manageable standards, so a challenge to Rule XI is nonjusticiable. Violations of the express oath, two-thirds, or Chief Justice requirements may remain reviewable.

\textbf{Reasoning.} ``Try'' has broad meanings; the Clause's three express requirements imply no others; ``sole'' means without interference. Impeachment is the Legislature's check on judges, review would threaten finality, and \textit{Powell v. McCormack} involved fixed textual limits. White found the claim justiciable---else the Senate could convict an official as ``a bad guy''---but upheld committee factfinding. Souter would intervene only if the Senate convicted ``upon a coin toss.''

\textbf{Examples.} \emph{Same-Side:} Senate limits on cross-examination. \emph{Opposite-Side:} a conviction by less than two-thirds. \emph{Fence-Sitter:} a summary conviction with no hearing.

\textbf{Critique.} Textualists praise the allocation of power and the refusal to let judges police the check on themselves; critics say Rehnquist overread ``sole'' and left ``try'' a parchment barrier, and that Souter's carve-out requires the merits judgment it disclaims.

\textbf{Key Quotation.} ``The conclusion that the use of the word `try' in the first sentence of the Impeachment Trial Clause lacks sufficient precision to afford any judicially manageable standard of review of the Senate's actions is fortified by the existence of the three very specific requirements that the Constitution does impose on the Senate when trying impeachments.'' --- Chief Justice Rehnquist, 506 U.S. at 230.

\end{document}
